% Колонтитул с указанием лицензии.
%
% Зачем: CC BY-SA 4.0, раздел 3(a), требует, чтобы сам материал отдавал лицензию и ссылку на неё.
% Комплект печатают и распространяют без README, поэтому указание лицензии должно быть на каждой странице.
%
% Почему без иконки CC:
% официальный значок CC просят не перерисовывать,
% а готовый пакет ccicons лежит в коллекции fontsextra — в CI это лишние полтора гигабайта ради одного значка.
% Текстовое обозначение лицензии со ссылками требование 3(a) выполняет полностью.
%
% Почему без имени:
% комплект перепечатывают и проводят разные организаторы, авторство остаётся в README и на странице релизов.
%
% Подключение:
% % Колонтитул с указанием лицензии.
%
% Зачем: CC BY-SA 4.0, раздел 3(a), требует, чтобы сам материал отдавал лицензию и ссылку на неё.
% Комплект печатают и распространяют без README, поэтому указание лицензии должно быть на каждой странице.
%
% Почему без иконки CC:
% официальный значок CC просят не перерисовывать,
% а готовый пакет ccicons лежит в коллекции fontsextra — в CI это лишние полтора гигабайта ради одного значка.
% Текстовое обозначение лицензии со ссылками требование 3(a) выполняет полностью.
%
% Почему без имени:
% комплект перепечатывают и проводят разные организаторы, авторство остаётся в README и на странице релизов.
%
% Подключение:
% % Колонтитул с указанием лицензии.
%
% Зачем: CC BY-SA 4.0, раздел 3(a), требует, чтобы сам материал отдавал лицензию и ссылку на неё.
% Комплект печатают и распространяют без README, поэтому указание лицензии должно быть на каждой странице.
%
% Почему без иконки CC:
% официальный значок CC просят не перерисовывать,
% а готовый пакет ccicons лежит в коллекции fontsextra — в CI это лишние полтора гигабайта ради одного значка.
% Текстовое обозначение лицензии со ссылками требование 3(a) выполняет полностью.
%
% Почему без имени:
% комплект перепечатывают и проводят разные организаторы, авторство остаётся в README и на странице релизов.
%
% Подключение:
% % Колонтитул с указанием лицензии.
%
% Зачем: CC BY-SA 4.0, раздел 3(a), требует, чтобы сам материал отдавал лицензию и ссылку на неё.
% Комплект печатают и распространяют без README, поэтому указание лицензии должно быть на каждой странице.
%
% Почему без иконки CC:
% официальный значок CC просят не перерисовывать,
% а готовый пакет ccicons лежит в коллекции fontsextra — в CI это лишние полтора гигабайта ради одного значка.
% Текстовое обозначение лицензии со ссылками требование 3(a) выполняет полностью.
%
% Почему без имени:
% комплект перепечатывают и проводят разные организаторы, авторство остаётся в README и на странице релизов.
%
% Подключение:
% \input{colophon.tex} в преамбуле,
% затем \pagestyle{ccfoot} (без номеров страниц) или \pagestyle{ccfootpage} (с номерами).

\RequirePackage{xcolor}
\RequirePackage{url}


% =====================================================================================================================
% Единая строка колонтитула: лицензия, ссылка на неё, ссылка на исходники.
%
% Шрифт здесь основной, а не \sffamily: guide.tex и rating.tex не задают \setsansfont,
% и polyglossia падает на латинском Modern Sans без кириллицы.
% Основной шрифт определён во всех документах.
%
% Ссылки написаны прямо в \url, а не через макросы:
% пакет url читает аргумент буквально и \url{\ccrepo} напечатал бы «\ccrepo».
\newcommand{\ccline}{%
  \hbox to \textwidth{%
    \color{black!45}\fontsize{7}{8.5}\selectfont
    \urlstyle{same}%
    CC BY-SA 4.0\enspace\url{https://creativecommons.org/licenses/by-sa/4.0/}%
    \hfil\url{https://github.com/askras/hexapawn-sama/}}%
  \vspace*{1pt}%
}

% =====================================================================================================================
% Тот же колонтитул с номером страницы справа.
% Под номер зарезервировано 2em, иначе \llap наезжает на ссылку на исходники.
\newcommand{\cclinepage}{%
  \hbox to \dimexpr\textwidth-2em\relax{%
    \color{black!45}\fontsize{7}{8.5}\selectfont
    \urlstyle{same}%
    CC BY-SA 4.0\enspace\url{https://creativecommons.org/licenses/by-sa/4.0/}%
    \hfil\url{https://github.com/askras/hexapawn-sama/}}%
  \hfil\llap{\color{black!45}\fontsize{7}{8.5}\selectfont\thepage}%
  \vspace*{1pt}%
}

% =====================================================================================================================
% Поля узкие (1,2 см), поэтому колонтитул ставим ближе к тексту, иначе он уходит за край листа.
\setlength{\footskip}{9mm}

% =====================================================================================================================
\makeatletter
% Вариант без номеров страниц: карточки, поле, таблица результатов.
\def\ps@ccfoot{%
  \let\@oddhead\@empty
  \let\@evenhead\@empty
  \def\@oddfoot{\ccline}%
  \def\@evenfoot{\ccline}%
}
% Вариант с номерами: методичка, правила, примеры.
\def\ps@ccfootpage{%
  \let\@oddhead\@empty
  \let\@evenhead\@empty
  \def\@oddfoot{\cclinepage}%
  \def\@evenfoot{\cclinepage}%
}
\makeatother
 в преамбуле,
% затем \pagestyle{ccfoot} (без номеров страниц) или \pagestyle{ccfootpage} (с номерами).

\RequirePackage{xcolor}
\RequirePackage{url}


% =====================================================================================================================
% Единая строка колонтитула: лицензия, ссылка на неё, ссылка на исходники.
%
% Шрифт здесь основной, а не \sffamily: guide.tex и rating.tex не задают \setsansfont,
% и polyglossia падает на латинском Modern Sans без кириллицы.
% Основной шрифт определён во всех документах.
%
% Ссылки написаны прямо в \url, а не через макросы:
% пакет url читает аргумент буквально и \url{\ccrepo} напечатал бы «\ccrepo».
\newcommand{\ccline}{%
  \hbox to \textwidth{%
    \color{black!45}\fontsize{7}{8.5}\selectfont
    \urlstyle{same}%
    CC BY-SA 4.0\enspace\url{https://creativecommons.org/licenses/by-sa/4.0/}%
    \hfil\url{https://github.com/askras/hexapawn-sama/}}%
  \vspace*{1pt}%
}

% =====================================================================================================================
% Тот же колонтитул с номером страницы справа.
% Под номер зарезервировано 2em, иначе \llap наезжает на ссылку на исходники.
\newcommand{\cclinepage}{%
  \hbox to \dimexpr\textwidth-2em\relax{%
    \color{black!45}\fontsize{7}{8.5}\selectfont
    \urlstyle{same}%
    CC BY-SA 4.0\enspace\url{https://creativecommons.org/licenses/by-sa/4.0/}%
    \hfil\url{https://github.com/askras/hexapawn-sama/}}%
  \hfil\llap{\color{black!45}\fontsize{7}{8.5}\selectfont\thepage}%
  \vspace*{1pt}%
}

% =====================================================================================================================
% Поля узкие (1,2 см), поэтому колонтитул ставим ближе к тексту, иначе он уходит за край листа.
\setlength{\footskip}{9mm}

% =====================================================================================================================
\makeatletter
% Вариант без номеров страниц: карточки, поле, таблица результатов.
\def\ps@ccfoot{%
  \let\@oddhead\@empty
  \let\@evenhead\@empty
  \def\@oddfoot{\ccline}%
  \def\@evenfoot{\ccline}%
}
% Вариант с номерами: методичка, правила, примеры.
\def\ps@ccfootpage{%
  \let\@oddhead\@empty
  \let\@evenhead\@empty
  \def\@oddfoot{\cclinepage}%
  \def\@evenfoot{\cclinepage}%
}
\makeatother
 в преамбуле,
% затем \pagestyle{ccfoot} (без номеров страниц) или \pagestyle{ccfootpage} (с номерами).

\RequirePackage{xcolor}
\RequirePackage{url}


% =====================================================================================================================
% Единая строка колонтитула: лицензия, ссылка на неё, ссылка на исходники.
%
% Шрифт здесь основной, а не \sffamily: guide.tex и rating.tex не задают \setsansfont,
% и polyglossia падает на латинском Modern Sans без кириллицы.
% Основной шрифт определён во всех документах.
%
% Ссылки написаны прямо в \url, а не через макросы:
% пакет url читает аргумент буквально и \url{\ccrepo} напечатал бы «\ccrepo».
\newcommand{\ccline}{%
  \hbox to \textwidth{%
    \color{black!45}\fontsize{7}{8.5}\selectfont
    \urlstyle{same}%
    CC BY-SA 4.0\enspace\url{https://creativecommons.org/licenses/by-sa/4.0/}%
    \hfil\url{https://github.com/askras/hexapawn-sama/}}%
  \vspace*{1pt}%
}

% =====================================================================================================================
% Тот же колонтитул с номером страницы справа.
% Под номер зарезервировано 2em, иначе \llap наезжает на ссылку на исходники.
\newcommand{\cclinepage}{%
  \hbox to \dimexpr\textwidth-2em\relax{%
    \color{black!45}\fontsize{7}{8.5}\selectfont
    \urlstyle{same}%
    CC BY-SA 4.0\enspace\url{https://creativecommons.org/licenses/by-sa/4.0/}%
    \hfil\url{https://github.com/askras/hexapawn-sama/}}%
  \hfil\llap{\color{black!45}\fontsize{7}{8.5}\selectfont\thepage}%
  \vspace*{1pt}%
}

% =====================================================================================================================
% Поля узкие (1,2 см), поэтому колонтитул ставим ближе к тексту, иначе он уходит за край листа.
\setlength{\footskip}{9mm}

% =====================================================================================================================
\makeatletter
% Вариант без номеров страниц: карточки, поле, таблица результатов.
\def\ps@ccfoot{%
  \let\@oddhead\@empty
  \let\@evenhead\@empty
  \def\@oddfoot{\ccline}%
  \def\@evenfoot{\ccline}%
}
% Вариант с номерами: методичка, правила, примеры.
\def\ps@ccfootpage{%
  \let\@oddhead\@empty
  \let\@evenhead\@empty
  \def\@oddfoot{\cclinepage}%
  \def\@evenfoot{\cclinepage}%
}
\makeatother
 в преамбуле,
% затем \pagestyle{ccfoot} (без номеров страниц) или \pagestyle{ccfootpage} (с номерами).

\RequirePackage{xcolor}
\RequirePackage{url}


% =====================================================================================================================
% Единая строка колонтитула: лицензия, ссылка на неё, ссылка на исходники.
%
% Шрифт здесь основной, а не \sffamily: guide.tex и rating.tex не задают \setsansfont,
% и polyglossia падает на латинском Modern Sans без кириллицы.
% Основной шрифт определён во всех документах.
%
% Ссылки написаны прямо в \url, а не через макросы:
% пакет url читает аргумент буквально и \url{\ccrepo} напечатал бы «\ccrepo».
\newcommand{\ccline}{%
  \hbox to \textwidth{%
    \color{black!45}\fontsize{7}{8.5}\selectfont
    \urlstyle{same}%
    CC BY-SA 4.0\enspace\url{https://creativecommons.org/licenses/by-sa/4.0/}%
    \hfil\url{https://github.com/askras/hexapawn-sama/}}%
  \vspace*{1pt}%
}

% =====================================================================================================================
% Тот же колонтитул с номером страницы справа.
% Под номер зарезервировано 2em, иначе \llap наезжает на ссылку на исходники.
\newcommand{\cclinepage}{%
  \hbox to \dimexpr\textwidth-2em\relax{%
    \color{black!45}\fontsize{7}{8.5}\selectfont
    \urlstyle{same}%
    CC BY-SA 4.0\enspace\url{https://creativecommons.org/licenses/by-sa/4.0/}%
    \hfil\url{https://github.com/askras/hexapawn-sama/}}%
  \hfil\llap{\color{black!45}\fontsize{7}{8.5}\selectfont\thepage}%
  \vspace*{1pt}%
}

% =====================================================================================================================
% Поля узкие (1,2 см), поэтому колонтитул ставим ближе к тексту, иначе он уходит за край листа.
\setlength{\footskip}{9mm}

% =====================================================================================================================
\makeatletter
% Вариант без номеров страниц: карточки, поле, таблица результатов.
\def\ps@ccfoot{%
  \let\@oddhead\@empty
  \let\@evenhead\@empty
  \def\@oddfoot{\ccline}%
  \def\@evenfoot{\ccline}%
}
% Вариант с номерами: методичка, правила, примеры.
\def\ps@ccfootpage{%
  \let\@oddhead\@empty
  \let\@evenhead\@empty
  \def\@oddfoot{\cclinepage}%
  \def\@evenfoot{\cclinepage}%
}
\makeatother
